\documentclass[12pt]{article}
\usepackage[utf8]{inputenc}
\usepackage[T1]{fontenc}
\usepackage{hyperref}
\title{Deko-Dusanga Language Resources for K-12 Curriculum Development}
\author{Compiled by Claude Sonnet 5 (claude-sonnet-5) for Andr\'as Kornai}
\date{}
\begin{document}
\maketitle

\section{Introduction}
Deko-Dusanga (ISO 639-3 uya; also Doko-Uyanga, Uyanga, Basanga, Dosanga, Iko) is a small Upper Cross River language of Cross River State, Nigeria. Speaker estimates range from about 200 (Wikipedia, Global Recordings Network) to a people group of 500 (Joshua Project). Little material exists. No school textbooks, primers or dictionaries were found in search results.

\section{Primary Educational Resources}
No school curriculum or primer was found.

\section{Linguistic and Reference Materials}
\begin{itemize}
\item Wikipedia, Uyanga language: \url{https://en.wikipedia.org/wiki/Uyanga_language}
\item Glottolog reference: Decker and Abraham 2022, a sociolinguistic profile that covers this language with two neighbouring languages: \url{https://glottolog.org/resource/reference/id/656699}
\item Endangered Languages Project bibliography: \url{https://endangeredlanguages.com/lang/3403/bibliography?hl=pt-br}
\end{itemize}

\section{Religious and Community Resources}
\begin{itemize}
\item Scripture Earth page for uya (lists scripture resources, including a link to a YouVersion Bible): \url{https://www.scriptureearth.org/00eng.php?iso=uya}
\item Joshua Project people group profile: \url{https://joshuaproject.net/people_groups/11592}
\end{itemize}

\section{Digital and Audio Resources}
\begin{itemize}
\item Global Recordings Network, Doko-Uyanga recordings: \url{https://globalrecordings.net/en/language/uya}
\end{itemize}

\section{Curriculum Development Context}
The language is not known to be taught in schools. Scripture and audio recordings are the main texts found. A curriculum would need new community-led documentation. The Decker and Abraham 2022 profile is a likely starting point for orthography and sociolinguistic context. I did not read it.
\end{document}